%% ============================================================
%% Chapter 2 — Methodology
%% ============================================================
%% chapter heading and short variant for headers/Table of Contents
\chapter[Description of the methodology (Name and surname of the chapter author in the case of a group thesis)]{Description of the methodology}\label{chapter:2_Methodology} 

This chapter should present a detailed description of the methods, tools, and techniques used in the thesis. It should develop the theoretical issues outlined in the introduction (see Chapter \ref{chapter:1_Introduction}) and clearly indicate how the author approached achieving the thesis objective.

First, the applied research or design methods should be described, together with their justification and the context of use. In empirical studies, the study group, sampling procedure, measurement instruments, and research procedures should be specified. For design or engineering theses, the technologies, programming environments, libraries, and frameworks should be presented, along with a brief discussion of the technical assumptions adopted.

Next, the project stages should be presented in a logical order - from requirements analysis and initial concepts, through implementation, to testing and validation of the results. It is advisable to supplement the description with block diagrams, UML diagrams, mathematical formulas, figures, or tables to clarify the outcomes.

\paragraph{Example structure of the methodology:} 
\begin{enumerate}
    \item Characteristics of methods and tools - a description of selected research techniques, algorithms, programming environments, or devices.
    \item Description of the implementation stages - actions presented in sequence, together with the justification of their use.
    \begin{enumerate}[label=\alph*)]
        \item Stage 1: requirements analysis.
        \item Stage 2: solution design.
        \begin{itemize}
            \item Stage 3: implementation and integration.
            \item Stage 4: testing and verification of results.
        \end{itemize}
    \end{enumerate}
    \item Summary - a brief overview of the adopted methodology and its relevance to achieving the thesis objectives.
\end{enumerate}

The entire content should be prepared coherently so that the reader can clearly follow the logic and sequence of actions leading to the thesis's results. Footnotes may be inserted using the command \verb|\footnote{Text.}|, e.g. \footnote{Their numbering follows the order of appearance in the thesis.}
